% 제출본 전용 — 부록으로 붙는다. 양식의 '3. 개발 및 운영 환경'.

\section{개발 및 운영 환경}

\subsection{환경 구성}

\begin{table}[htbp]
  \centering
  \caption{개발 및 운영 환경}
  \label{tab:v2-env}
  \begin{tabular}{lp{10.0cm}}
    \toprule
    \textbf{구분} & \textbf{내용} \\
    \midrule
    형태 & \textbf{단독 PC.} 서버·클라우드·데이터베이스를 쓰지 않는다.
      결과는 파일(\texttt{parquet} · \texttt{json} · \texttt{csv})로 남는다 \\
    운영체제 & Windows 11 Pro (22621). 경로 처리 외에 OS 의존은 없다 \\
    언어 & Python 3.13.7 \\
    주요 패키지 & pandas · numpy · statsmodels · scipy · matplotlib ·
      shapely (법정동 경계) · PyMuPDF (공고 PDF 판독) · python-pptx (판넬) \\
    외부 자료 & 국토교통부 실거래 상세 자료 오픈 API\citep{molit_api} ·
      서울시보 · 전자관보 공고 원문 \\
    문서 & \LaTeX\ (XeLaTeX + kotex) · Markdown $\to$ HTML \\
    공개 & GitHub 저장소\citep{repo} · GitHub Pages\citep{site} \\
    \bottomrule
  \end{tabular}

  \begin{minipage}{\linewidth}\footnotesize
  \textit{주.} 별도의 기자재나 재료를 쓰지 않았으므로 기자재 목록과 재료비
  항목은 해당 사항이 없다. 추가 비용은 0원이다(10.5절).
  \end{minipage}
\end{table}

\textbf{데이터베이스를 쓰지 않은 것은 설계 판단이다.} 자료가 47만 행으로
단일 PC 메모리에 올라가고, 분석이 질의가 아니라 \textbf{전체 재계산}이기
때문이다. 파일로 두면 결과가 바이트 단위로 비교되므로 재현 확인이 간단해진다
--- 실제로 클린 환경 재현에서 결과 파일 25개를 해시로 대조했다(3.6절).

\subsection{의존 목록 관리}

\texttt{requirements.txt}는 \texttt{scripts/check\_deps.py}가 지킨다. 저장소의
모든 Python 파일을 AST 로 훑어 최상위 import 를 모으고, 목록과 어긋나면 알린다.
함수 안의 늦은 import 까지 본다.

이 장치를 만든 이유가 있다. 2026-10-15 클린 환경 재현에서 \texttt{shapely}와
\texttt{markdown}이 목록에서 빠져 있던 것이 드러났다. 작업 환경에는 다른 경로로
깔려 있어 몇 주 동안 보이지 않았고, \textbf{새 가상환경에서만 12단계에서
멈췄다.} 받는 사람만 겪고 만드는 사람은 못 보는 고장이다.

\subsection{환경 관련 문제와 대응}

\textbf{윈도의 경로 길이 한계(260자).} 깊은 폴더에서 \texttt{pip install}을 하면
\texttt{statsmodels}의 시험 자료 경로가 한계를 넘어 \textbf{그 패키지만 빠진 채
설치가 끝난다.} 137자짜리 폴더에서 실제로 겪었다(파일 경로 266자). 짧은 경로에
두거나 긴 경로 지원을 켠다.

\textbf{콘솔 인코딩.} 한글을 찍는 스크립트는 \texttt{lib.io.setup\_stdout()}으로
표준출력을 UTF-8 로 돌린다. 그러지 않으면 cp949 콘솔에서 죽는다.

\textbf{저장소 용량.} 그림을 매번 커밋해 \texttt{.git}이 936\,MB 다. 이력을 다시
쓰는 일이라 이번 학기에는 고치지 않았다(10.4절).
